\section*{Limitations}
\label{sec:limitations}

We discuss potential limitations of the method, the data, and the tooling, and
the design choices we use to minimize each.

\paragraph{The structured prior is synthesized by a language model.}
Both the per-app FSMs and the category libraries $L_C$ are synthesized by a
frozen Claude Opus model from a small bundle of trajectories, which in principle
risks two failure modes: leakage of app-specific tokens into the supposedly
transferable upper layer, and low-quality or spurious archetypes. We minimize
both structurally rather than by stylistic discipline. A linter
(\texttt{lint\_layer2}) rejects any LAYER~2 block---and any diff proposed during
online evolution---that contains an app name or a source-pool resource
identifier, so the upper layer's transferability is enforced at every point where
it could be violated, and the synthesized libraries are additionally audited
before use.

\paragraph{The transferable library can have coverage gaps.}
A target app may exercise functionality that no same-category source app covers,
so a category library is not guaranteed to contain a directly applicable
archetype. We mitigate this in three ways: a three-tier resolution at injection
time (app-level FSM $\to$ category $L_C$ $\to$ a generic bootstrap), so a missing
tier falls back rather than fails; entry-level retrieval in the cross-benchmark
setting, which injects only the few archetypes a retriever judges relevant to the
task instead of the whole library; and a far-transfer panel that receives no
injection at all, degrading byte-identically to the zero-shot prompt so that the
absence of a matching prior can never harm the agent.

\paragraph{The cross-benchmark target has no seed dimension.}
MobileWorld tasks are hard-coded single instances with no random seed, which
reduces the instance diversity available during adaptation and could in principle
let an evolving FSM fit the specific adaptation instances. We minimize this with a
strictly task-disjoint adaptation/evaluation split, so evaluation is always on
tasks unseen during adaptation; by selecting diverse and multi-application
adaptation tasks so each app receives varied adaptation coverage; and by relying
on symbolic evolution, which distills generalizable workflow patterns from
observed failures rather than memorizing trajectories and is therefore
comparatively insensitive to the missing seed axis.

\paragraph{App-specific structure is rediscovered in context.}
Our implementation collapses the offline extraction of a target app's LAYER~1
into an in-context rediscovery during the adaptation loop, rather than building a
target LAYER~1 ahead of time. The strong visual encoder makes this viable, and
our cross-benchmark analysis shows that the alternative---transplanting a
source-environment LAYER~1 into the prompt---is in fact counterproductive, since
source-grounded state descriptions can displace the agent's own grounded
behavior. Declining to transplant app-specific structure is therefore aligned
with what the data shows, not a shortcut.

\paragraph{Environment and tooling instrumentation.}
Online mobile-agent benchmarks occasionally surface instrumentation gaps---for
instance a virtual-table extension missing from the host database build, or a
natural-language app name the launcher cannot resolve---that can make individual
tasks error out independently of agent behavior. We minimize their effect by
detecting and fixing those we can at the infrastructure level (for example,
restoring the required database extension and the media directory the audio tasks
expect, and provisioning fresh, state-isolated emulator containers for every run)
and by reporting any residual environment-induced failures uniformly across all
baselines, so they cancel in every comparison rather than favoring one arm.
